Call a pairwise orthogonal subset $F$ of an orthoset an
\emph{admissible figure} when
\[
 \cl F\cap N(A)=\cl(F\setminus A)\qquad(A\subseteq F).
\]
The quantifier over figures is important: their supports need not
represent every orthogonal family in the complete lattice.

\begin{theorem}[Figure complements and finite flags]\label{thm:flags}
For any orthoset $E$, the following are equivalent: every pairwise
orthogonal subset is admissible; and, for every maximal pairwise
orthogonal subset $D$ and distinct $u,v\in E$,
\[
 D\subseteq N(u)\cup N(v)\quad\Longrightarrow\quad u\perp v.
\]
The equivalence includes infinite subsets, with the axiom of choice.
For every poset $L$, each finite pairwise orthogonal subset of $M(L)$
is admissible. This finite assertion does not imply orthomodularity
of $\Sigma(M(L))$, and it does not extend to all infinite figures.
\end{theorem}
For finite graphs the displayed condition is Sumner's definition
\cite[p.~492]{Sumner}. It is the normality condition of
Paseka--Vetterlein in their finite-rank setting
\cite[Definition~2.4 and Lemma~2.5, pp.~5--6]{PV}.
It is weaker than the basis criterion for an orthomodular support
lattice. Section~\ref{sec:figure-comparisons} gives explicit finite
and infinite witnesses.

For subsets $A,B$ of an orthoset, $A\perp B$ means $A\subseteq N(B)$.
The object and arrow conventions also matter. A complete
ortholattice is a complete lattice with an antitone involution $C$
satisfying $x\wedge Cx=0$ and $x\vee Cx=1$. For orthosets $E,F$,
a \emph{saturated correspondence} is a map
$\alpha:E\to\Sigma(F)$ such that, for every $A\subseteq E$,
\[
 x\in\cl_E A\quad\Longrightarrow\quad
 \alpha(x)\subseteq\bigvee_{a\in A}\alpha(a).
\]

\begin{theorem}[The complete support correspondence]\label{thm:correspondence}
Every complete ortholattice $K$ is represented by the orthoset
$K\setminus\{0\}$, with $x\perp y$ if and only if $x\le Cy$:
\[
 K\xrightarrow{\sim}\Sigma(K\setminus\{0\}),\qquad
 s\longmapsto\{x\ne0:x\le s\}.
\]
For every pair of orthosets $E,F$, saturated correspondences are
in bijection with maps $f:\Sigma(E)\to\Sigma(F)$ preserving all
joins, including the empty join. The inverse formulas are
\[
 f(S)=\bigvee_{x\in S}\alpha(x),\qquad
 \alpha(x)=f(\cl_E\{x\}).
\]
The identity is $x\mapsto\cl_E\{x\}$ and composition is
$(\beta\alpha)(x)=\bigvee_{y\in\alpha(x)}\beta(y)$.
Consequently these two categories are equivalent. If $g$ is the
order right adjoint of $f$, the associated dagger is
$f^\dagger=C_EgC_F$; it is characterized by
\[
 f(S)\perp T\quad\Longleftrightarrow\quad S\perp f^\dagger(T).
\]
Ordinary orthogonality-preserving functions are not the arrows in
this equivalence.
\end{theorem}
The object axioms are exactly the Spat axioms on archive
sheets~50--52 of cote~156-9 \cite{Groth}. We specify the arrow
condition explicitly so that this statement does not depend on an
unspecified meaning of morphism.
